\section{Introduction}

Ce chapitre a pour but de présenter d'une manière générale le projet réalisé dans le cadre de ce stage de fin d'études en quatre parties :

\begin{itemize}
    \item La première partie concerne l'organisme d'accueil, TAMTAM International, avec plus de détails sur ses fondateurs, ses clients ainsi que les solutions logicielles qu'elle propose.
    
    \item La deuxième partie est dédiée à la présentation des formations dispensées en début de stage, expliquant comment elles se sont déroulées.
    
    \item La troisième partie présente le contexte général du stage.
    
    \item Enfin, la quatrième partie décrit la méthodologie de travail adoptée tout au long du projet.
end{itemize}

\section{Contexte Général}

\subsection{Présentation de l'entreprise Tamtam International}

Tamtam International est une entreprise belge innovante, fondée en 2012 à Solvay, experte en développement et l'intégration de solutions web et mobiles. Avec l'ouverture d'un bureau à Marrakech en 2015, cela a renforcé son positionnement, elle est devenue intervenant incontournable dans le domaine des solutions informatiques destinées aux professionnels de l'économie.

La philosophie de Tamtam International repose sur le développement d'un écosystème dédié à la gestion de communautés, a pour objectif de devenir le leader du marché dans ce domaine. L'entreprise collabore avec une clientèle diversifiée, allant des experts-comptables et cabinets fiduciaires comme le Cabinet Expert-comptable Belgique Deg \& Partner, et se consacre à élaborer un avenir numérique plus sûr pour ses partenaires.

\subsection{Fondateur Tamtam International}

\textbf{Emmanuel Dégrève} est fondateur et CEO de Tamtam International, une référence dans le domaine fiscal en Belgique. Expert en matière d'impôt des personnes physiques (IPP) et des processus de création de sociétés. En outre, il préside le Forum For the Future, une fondation qui organise des événements annuels pour les professions économiques, tels que le congrès à Brussels Expo. Son expertise approfondie, ainsi que son leadership entrepreneurial font de lui une figure influente du monde financier et de la gestion.

\subsection{Activité de Tamtam International}

L'entreprise Tamtam se distingue des autres entreprises par sa philosophie, elle ne se spécialise pas dans le développement des solutions informatiques pour d'autres clients à la demande, mais elle se focalise sur le développement d'un écosystème pour la gestion de communautés et vise à devenir le leader sur le marché.

\subsection{Fiche technique de Tamtam International}

\begin{table}[H]
\centering
\begin{tabular}{|l|l|}
\hline
\textbf{Raison sociale} & Tamtam International \\
\hline
\textbf{CTO} & Emmanuel Dergève \\
\hline
\textbf{CEO} & Bruno de l'Escaille \\
\hline
\textbf{Adresse} & Résidence El Batoul, Rue Mohamed El Bequal Bloc B, Appt N°20, Marrakech – Médina \\
\hline
\textbf{RC} & 53375 (Tribunal de Marrakech) \\
\hline
\textbf{Forme juridique} & SRLAU \\
\hline
\textbf{Capital} & 99 000 DH \\
\hline
\textbf{Année de création} & 2013 \\
\hline
\textbf{Email} & tamtampro.it@gmail.com \\
\hline
\textbf{Téléphone} & +212 (0) 524 156 312 \\
\hline
\textbf{Taille} & 11-50 employés \\
\hline
\end{tabular}
\caption{Fiche technique de Tamtam International}
\end{table}

\subsection{Clients de Tamtam International}

Tamtam vise le marché des fiduciaires et des expert-comptables en Belgique. Ses principaux clients sont :

\begin{itemize}
    \item \textbf{Deg Partners} : un cabinet comptable spécialisé dans l'accompagnement des professions libérales, des artistes et des associations (ASBL). Il se distingue par son expertise en matière de fiscalité et de transition vers le statut d'entreprise.
    
    \item \textbf{IEC} : (Institut des Experts-comptables et des Conseils fiscaux – Belgique) un organisme professionnel regroupant les experts-comptables et les conseillers fiscaux en Belgique.
    
    \item \textbf{UP} : (UnifiedPost Group) Un fournisseur de premier plan de solutions innovantes pour des opérations commerciales fluides.
    
    \item \textbf{FFF} : (Forum For the Future) une fondation qui organise un événement annuel rassemblant les professionnels des secteurs économiques.
    
    \item \textbf{ITAA} : (Institute for Tax Advisors Accountants) institut belge qui représente les conseillers fiscaux et les comptables.
end{itemize}

\subsection{Produits de TamTam International}

La société Tamtam s'engage à la délivrance des divers produits qui servent à la fois dans la gestion événementielle, la satisfaction des besoins des professionnels économiques et experts-comptables. Principalement, Tamtam travaille sur un écosystème composé de :

\begin{itemize}
    \item \textbf{API Tamtam} : une API développée en Symfony qui regroupe les modules.
    
    \item \textbf{Home} : l'application centrale qui représente l'interface d'authentification.
    
    \item \textbf{Emailling} : permet aux clients d'envoyer des emails à des contacts.
    
    \item \textbf{Blog} : offrant un espace personnel aux utilisateurs.
    
    \item \textbf{Event} : application web dédiée à la gestion des événements.
    
    \item \textbf{Talk} : application qui facilite la communication entre collègues.
    
    \item \textbf{WhereToBill} : un moteur de routage pour le transfert des documents.
    
    \item \textbf{SMS} : application pour informer les clients par SMS.
    
    \item \textbf{Billing} : permet la gestion du paiement en ligne.
    
    \item \textbf{PowerTeam} : plateforme dédiée à la gestion des activités des fiduciaires.
    
    \item \textbf{Webinar} : plateforme offre aux orateurs un partage de leurs formations.
    
    \item \textbf{Dashboard} : application contient des tableaux de bord personnalisés.
end{itemize}

\section{Présentation de la formation}

Afin de permettre une intégration fluide de stagiaires dans le processus de travail au sein de Tamtam International, les responsables des stagiaires ont organisé une formation qui sert à mettre en point les prérequis techniques nécessaires pour commencer le travail.

\subsection{Installation de l'environnement de développement}

Tout en respectant un ensemble de spécifications, nous avons installé et configuré l'environnement TTP sur nos machines :

\begin{itemize}
    \item Installation du système d'exploitation UNIX/LINUX (Ubuntu).
    \item Installation de PHP 7.4, PHP-FPM ainsi toutes ses dépendances.
    \item Installation de MySQL 5.7.
    \item Installation et configuration de NGINX à la place d'Apache.
    \item Configuration des projets (Composer, Node.js, Yarn, Git).
    \item Import de la base de données.
    \item Installation du projet APITTP et ses dépendances.
    \item Installation du projet HOME pour l'authentification.
end{itemize}

\subsection{Comprendre l'environnement technique}

Après avoir terminé la configuration de notre environnement de travail en local, nous avons bénéficié de deux formations techniques :

\begin{itemize}
    \item Une formation sur \textbf{Symfony}, axée sur la compréhension des composants du framework.
    \item Une formation sur \textbf{React} et \textbf{Redux}, permettant d'aborder les concepts fondamentaux du développement frontend.
    \item Introduction à la méthodologie Agile Scrum.
    \item Présentation du processus DevOps utilisé pour le déploiement.
end{itemize}

\subsection{Intégration des stagiaires}

Une fois que nous avons terminé la période de formation, chacun de nous est intégré dans un projet. Le chef du projet nous a donné une présentation détaillée du projet et de son fonctionnement.

\section{Présentation du sujet de stage}

\subsection{Contexte général du projet}

Tamtam International, basée en Belgique, est progressivement devenue un acteur majeur grâce aux solutions qu'elle propose dans le domaine de la fiscalité. Son succès s'appuie sur sa capacité à concevoir des projets ciblés, conçus en étroite collaboration avec les utilisateurs finaux.

Fort du succès rencontré par sa suite d'outils numériques, Tamtam a choisi de consolider sa présence sur le marché en développant un écosystème d'applications sous l'intitulé Community Management.

\subsection{Description du projet}

\textbf{PowerTeam} est une solution avancée qui offre une plateforme conçue spécifiquement pour les fiduciaires belges pour les accompagner dans la gestion budgétaire des dossiers clients et le suivi de leur activité.

Bien que PowerTeam soit déjà en production, son déploiement a mis en évidence plusieurs améliorations afin de répondre aux exigences des clients.

\subsection{Problématique}

Dans le cadre de la transition numérique, les entreprises cherchent à améliorer leurs processus internes pour gagner en efficacité et en qualité de service. Dans les cabinets comptables ou services financiers, la gestion des prestations telles que les déclarations fiscales, les bilans ou les audits repose fréquemment sur des outils souvent obsolètes.

Face à ces constats, il devient essentiel de repenser l'architecture et les fonctionnalités de ces outils de gestion intégrée des prestations.

\section{Méthodologie Adoptée}

\subsection{Méthodologie Scrum}

En termes de méthodes de travail, l'équipes de \textit{Tamtam} s'efforcent de suivre un ensemble de principes jugés pragmatiques pour développer et gérer les différents projets en cours. Ces principes découlent des méthodes Agiles.

Pour \textit{Tamtam}, la méthode agile la plus répandue est Scrum. Scrum est un Framework méthodologique qui facilite :

\begin{itemize}
    \item Flexibilité et réactivité
    \item Une grande capacité d'adaptation au changement grâce à des cycles courts d'itération.
    \item Une approche organisée qui permet de tester, valider et ajuster chaque étape du projet.
end{itemize}

\subsection{Les rôles Scrum}

La méthode Scrum définit 3 rôles :

\begin{description}
    \item[Product Owner] détient l'expertise fonctionnelle du projet et celui qui prend les décisions essentielles à la priorisation des développements. Dans Tamtam, le product Owner est M. Emmanuel Degrève.
    
    \item[Scrum Master] celui qui optimise la capacité de production de l'équipe. Dans notre équipe, le Scrum Master est Mme. EL JEMLI Safae.
    
    \item[Équipe de développement] composée de professionnels qualifiés, prend en charge de concevoir, développer et tester les fonctionnalités du produit.
end{description}

\subsection{Gestion des tâches et suivi}

Pour la gestion des sprints, Tamtam International a opté pour la solution Notion pour le suivi des tâches. Les colonnes de gestion sont :

\begin{itemize}
    \item \textbf{Backlog} : Liste de toutes les tâches potentielles, non encore planifiées.
    \item \textbf{To Do} : Tâches sélectionnées pour être traitées durant le sprint.
    \item \textbf{Urgent} : Tâches prioritaires à traiter en premier.
    \item \textbf{In Progress} : Tâches actuellement en cours de développement.
    \item \textbf{Pull Request} : Tâches dont le code a été soumis en Pull Request sur GitHub.
    \item \textbf{Preprod} : Tâches déployées en préproduction pour les tests finaux.
    \item \textbf{Validation} : Tâches en attente de validation ou d'approbation finale.
    \item \textbf{Done} : Tâches complètement terminées, validées et mises en production.
end{itemize}

\section{Conclusion}

Dans ce chapitre nous avons mis le point sur la présentation de l'entreprise tout en présentant ses produits et sa fiche technique. Après, nous avons précisé le contexte général du projet, en fixant ses axes majeurs et en détaillant chacun. Enfin nous avons présenté la méthode de travail suivi pour la réalisation du projet.
